\section{The main theorem}
\label{ch:main-theorem}

The statement below is the formalization target of the whole project. It matches,
definition for definition, the Lean file
\texttt{MIPRE/HaltingGameValue.lean}, which states it on top of Mathlib alone (theorem
\texttt{halting\_reduces\_to\_gameValue}): Turing machines are Mathlib's
\texttt{Nat.Partrec.Code}, games are game descriptions
(Definition~\ref{def:game-description}), and the value is the synchronous value
(Definition~\ref{def:sync-value}).

\begin{theorem}[Halting reduces to the game value; Theorem 12.2 of~\cite{JNVWY20},
synchronous form]
\label{thm:main}
\uses{thm:halting,def:game-description,def:sync-value,def:re}
\lean{HaltingGameValue.halting_reduces_to_gameValue}
\leanok
There is a computable map $g$ from Turing machines to game descriptions such that for
every Turing machine $\machine$:
\begin{enumerate}
\item (\textbf{Completeness}) if $\machine$ halts on the empty input, then
  $\synval(\game_{g(\machine)}) = 1$;
\item (\textbf{Soundness}) if $\machine$ does not halt on the empty input, then
  $\synval(\game_{g(\machine)}) \leq \tfrac12$.
\end{enumerate}
\end{theorem}

Theorem~\ref{thm:main} follows from Theorem~\ref{thm:halting} by converting the
normal form verifier at a fixed index into an explicit game description (a finite
object obtained by running the sampler and decider), and relaxing polynomial-time
computability of the reduction to computability, per Section~\ref{sec:departures}. A game
description rejects unequal answers to equal questions by fiat
(Definition~\ref{def:game-description}), which costs nothing here: against a tracial state
a projective family is perfectly consistent, so the terms the fiat deletes carry no weight
in $\synval$, which is unchanged, and it can only lower $\valstar$ --- the safe direction
for the soundness clause, while the completeness witness is synchronous and survives.
The two passages this needs from Theorem~\ref{thm:halting}, whose Lean assembly concludes in
$\valstar$, are formalized: the transport of a value-$1$ PCC strategy along the relabelings of
the tabulation (\texttt{MIPRE.Verifier.gameValue\_toGame\_eq\_one\_doubled}, through
\texttt{MIPRE.syncValue\_eq\_of\_equiv}) for completeness, and
\texttt{MIPRE.Verifier.gameValue\_toGame\_le\_of\_valStar\_le\_doubled} for soundness; both are
stated directly in \texttt{HaltingGameValue.gameValue}, by
Remark~\ref{rem:standalone-restatement}'s identification of the two vocabularies, and both are
read against the \emph{doubled} game of $\verifier_n$, which is how the tabulation matches it
without asking the verifier to be synchronous (Remark~\ref{rem:compression-abstract}, item 3).
On the tabulation the completeness bridge is \texttt{MIPRE.Halting.gameValue\_tab\_eq\_one}.
What the proof of Theorem~\ref{thm:main} still waits on is the obligations of
Theorem~\ref{thm:halting} that remain open, and mathematically that is the compressor alone:
the two distinguished strings and the semidecider are theorems, kept as fields of
\texttt{MIPRE.Halting.Obligations} until they are assembled into an inhabitant. The
tabulation is not among them at all, and neither is the record of its agreement with
$\verifier_n$ in the form the completeness bridge consumes: that record is
\texttt{MIPRE.Halting.tab\_match} (Remark~\ref{rem:compression-abstract}).

\begin{remark}[The standalone restatement]
\label{rem:standalone-restatement}
\uses{def:game-description,def:sync-value}
\lean{HaltingGameValue.gameValue, HaltingGameValue.GameData.equivTuple,
  HaltingGameValue.GameData.syncGame, HaltingGameValue.GameData.game,
  HaltingGameValue.GameData.gameValue_eq_syncValue,
  HaltingGameValue.GameData.gameValue_le_quantumValue}
\leanok
The Lean main statement is deliberately self-contained: the file
\texttt{MIPRE/HaltingGameValue.lean} imports nothing from the rest of this development,
re-declaring the synchronous game, the synchronous strategy and the value on top of
Mathlib alone, so that Theorem~\ref{thm:main} can be read without any of the vocabulary
built in Section~\ref{ch:foundations}. In that reduced vocabulary \texttt{gameValue} is
the synchronous value of Definition~\ref{def:sync-value} and \texttt{equivTuple}
exhibits the canonical G\"odel numbering asserted in
Definition~\ref{def:game-description} as an explicit equivalence between game
descriptions and a tuple of integers and lists, which is what supplies the
\texttt{Encodable} instance the computability clause quantifies over. The duplication is
the price of a statement that stands alone, and keeping the two vocabularies in step is a
maintenance obligation rather than a mathematical one: the definitions match one for one,
and this is the list of the ones that do. The correspondence is not left to inspection where
it is used: \texttt{GameData.syncGame} and \texttt{GameData.game} read a description in the
vocabulary of Section~\ref{ch:foundations} with the same distribution and predicate, and
\texttt{GameData.gameValue\_eq\_syncValue} proves the two synchronous values equal --- the
two structures of synchronous strategy, one carrying \texttt{POVM}s with a projectivity
certificate and the other a \texttt{ProjectiveMeasurement}, repackage into each other with
the same dimension and the same operators, so the two suprema are over the same set of
numbers. \texttt{GameData.gameValue\_le\_quantumValue} is the comparison
$\synval \leq \valstar$ (Lemma~\ref{lem:sync-le-valstar}) in that vocabulary.
\end{remark}

\begin{corollary}[Halting reduces to the quantum value]
\label{cor:main-quantum}
\uses{thm:halting,thm:main,lem:sync-le-valstar,def:q-value,def:game-description}
With $g$ as in Theorem~\ref{thm:main}: if $\machine$ halts on the empty input then
$\valstar(\game_{g(\machine)}) = 1$, and if $\machine$ does not halt on the empty input
then $\valstar(\game_{g(\machine)}) \leq \tfrac12$.
\end{corollary}

Both forms come from the two items of Theorem~\ref{thm:halting} directly: the value-$1$
witness is a PCC, hence synchronous, strategy, so it gives value $1$ in both senses; and
the soundness conclusion is $\valstar(\game_\machine) \leq \frac12$, which gives
$\synval(\game_\machine) \leq \frac12$ by Lemma~\ref{lem:sync-le-valstar}. The quantum
form is the one Definition~\ref{def:mipstar} asks for, and it is the primary one: no
transport from $\synval$ to $\valstar$ is involved anywhere, the pipeline's soundness being
carried in $\valstar$ at every step (Remark~\ref{rem:bipartite-route}). The synchronous
reading of Theorem~\ref{thm:main} is the derived one, obtained from this corollary by the
free comparison $\synval \leq \valstar$.

\begin{corollary}[Uncomputability of the game value]
\label{cor:value-uncomputable}
\ledgernode{1.7}
\uses{thm:main,cor:main-quantum,thm:halting-undecidable}
There is no algorithm that, given a game description $g$ with the promise that
$\synval(\game_g) = 1$ or $\synval(\game_g) \leq \frac12$, decides which is the
case; the same holds with $\valstar$ in place of $\synval$
(Corollary~\ref{cor:main-quantum}).
\end{corollary}

\begin{lemma}[$\MIPstar \subseteq \RE$]
\label{lem:mipstar-sub-re}
\uses{def:mipstar,def:re,lem:value-lower-approx}
\lean{MIPRE.MIPStar.isRE, MIPRE.MIPStar.exists_semidecider}
\leanok
Every language in $\MIPstar$ is recursively enumerable.
\end{lemma}

\begin{proof}
\leanok
Let $z \mapsto \game_z$ witness $L \in \MIPstar$. Under the promise, $z \in L$ if and only
if $\valstar(\game_z) > \frac12$: the value is $1$ on $L$ and at most $\frac12$ off it. By
Lemma~\ref{lem:value-lower-approx} the set of game descriptions $g$ with
$\valstar(\game_g) > \frac12$ is recursively enumerable, and $L$ is its preimage under the
computable map $z \mapsto \game_z$; the preimage of an r.e.\ set under a computable map is
r.e. Concretely: given $z$, compute $\game_z$, enumerate tensor-product strategies and accept
upon finding one of value greater than $\frac12$; this semi-decides $L$. In Lean the class
is the computable one (Definition~\ref{def:mipstar}), and
\texttt{MIPStar.exists\_semidecider} gives the semidecider as a program of the ambient
model.
\end{proof}

\begin{theorem}[$\MIPstar = \RE$]
\label{thm:mipstar-eq-re}
\ledgernode{1}
\uses{def:mipstar,def:re,thm:halting,cor:main-quantum,thm:halting-undecidable,lem:value-lower-approx,lem:mipstar-sub-re}
$\MIPstar = \RE$.
\end{theorem}

For $\RE \subseteq \MIPstar$: the halting problem is $\RE$-complete
(Theorem~\ref{thm:halting-undecidable}) and Theorem~\ref{thm:halting} provides a
polynomial-time verifier for it, whose completeness--soundness gap is already in the form
Definition~\ref{def:mipstar} requires (Corollary~\ref{cor:main-quantum}); no transport
between the two values is needed. $\MIPstar \subseteq \RE$ is
Lemma~\ref{lem:mipstar-sub-re}.

\begin{remark}
The polynomial-time version of Theorem~\ref{thm:main} (Theorem 12.2
of~\cite{JNVWY20} as stated there) is recorded by Theorem~\ref{thm:halting}; the
computable version above is what the Lean main statement asserts, and is sufficient
for every consequence in Section~\ref{ch:downstream}.
\end{remark}
